\chapter{Introducción}\label{cap:Introduccion}
% Índice v4, cap. 1: el reto, el cambio de paradigma y la solución en 2-3 páginas,
% sin detalle técnico.
%  - El reto: el agricultor debe dedicarse a su explotación, pero la normativa le
%    obliga a registrar tratamientos, fertilizaciones y labores.
%  - El cambio de paradigma: los agentes de IA escriben buena parte del código; el
%    valor del ingeniero se desplaza a especificar, decidir y verificar (SDD).
%  - La solución: aplicación generada con Claude Code en VS Code a partir de una
%    especificación, con RAG que cita la fuente y pipeline CI/CD.
\pendiente{introducción: el reto, el cambio de paradigma y la solución}

\section{Motivación}\label{sec:Motivacion}
% Gira en torno al agricultor, no a la tecnología. Cuatro pasos:
%  1. El problema real: registros obligatorios y normativa (PAC, fitosanitarios,
%     nutrición de suelos, cuaderno digital y SIEX, RD 1054/2022). Comprobar el
%     calendario vigente y citar la norma.
%  2. La fricción: cuantificarla con fuentes (MAPA, organizaciones agrarias,
%     consejería de CLM) y con 2-3 entrevistas. No inventar cifras.
%  3. La propuesta: lenguaje natural por voz o texto, avisos de plazos y dudas
%     normativas resueltas citando la fuente (justifica el RAG).
%  4. Por qué SDD y agentes: requisitos legales en una especificación versionada y
%     trazados hasta su prueba.
\pendiente{motivación en cuatro pasos}

\section{Objetivos}\label{sec:Objetivos}
% Objetivo general y seis objetivos específicos numerados; el capítulo 6 los
% revisa uno a uno.
\pendiente{objetivo general y objetivos específicos}

\section{Estructura de la memoria}\label{sec:Estructura}
% Un párrafo con el contenido de cada capítulo.
\pendiente{un párrafo por capítulo}
